\documentclass[10pt,a4paper]{amsart}
\usepackage[margin=19mm]{geometry}
\usepackage{amsmath,amssymb,mathtools}
\usepackage[scr=rsfs,cal=euler]{mathalfa}
\usepackage{xfrac}
\usepackage[unicode,hidelinks,bookmarks=false]{hyperref}
\newcommand{\ie}{i.e.\ }
\newcommand{\cf}{cf.\ }
\newcommand{\resp}{respectively\ }
\newcommand{\Subsection}{\subsection}
\providecommand{\nobreakdash}{\nobreak\hbox{-}}
\setlength{\emergencystretch}{2em}
\multlinegap0pt
\usepackage{bm}
\usepackage{bbm}
\usepackage{dsfont}
\usepackage{xspace}
\usepackage{cancel}
\usepackage{mathtools}
\usepackage{cleveref}
\usepackage{amsthm}
\makeatletter
\@ifundefined{theorem}{\newtheorem{theorem}{Theorem}}{}
\@ifundefined{lemma}{\newtheorem{lemma}[theorem]{Lemma}}{}
\makeatother
\title[Discrete Curvatures and Convex Polytopes]{Discrete Curvatures and Convex Polytopes}
\author{Jes\'us A. De Loera, Jillian Eddy, Sawyer Jack Robertson, Jos\'e Alejandro Samper}
\date{Source: arXiv:2510.11894v1 (CC BY 4.0)}
\begin{document}
\maketitle

\noindent\emph{Source and licence.} Discrete Curvatures and Convex Polytopes. Version arXiv:2510.11894v1 (CC BY 4.0), distributed under the Creative Commons Attribution 4.0 International licence, as recorded in the arXiv licence metadata for this exact version and shown on the abstract page https://arxiv.org/abs/2510.11894v1. Full rights notice: licenses/arxiv-2510.11894-CC-BY-4.0.txt.\par\smallskip

\noindent\textit{Editorial scope.} Counted unit: the Lemma whose source label is \texttt{lem:block-diag} in the article (source0.tex lines 2352--2393, source label \texttt{lem:block-diag}), delivered together with the article's own complete original proof of it (source0.tex lines 2395--2606, one \texttt{proof} environment of 212 lines). No mathematics is written, repaired or re-derived: statement and proof are the article's own text line for line. Required context retained uncounted: none. Every definition, display and label of those lines is retained unchanged. Omitted from this package and not redistributed: 1--2351, 2394--2394, 2607--2955. Nothing inside a retained excerpt is omitted or altered: every retained body is delivered exactly as the article's lines are written, so no omission receipt of any kind accompanies this package. Citations: the retained excerpts carry no citation command at all, so no bibliography entry of the article is transcribed and no reference is redistributed. Reference adaptation: none was needed and none was made; every reference inside the delivered unit resolves inside it, so no unresolved reference or citation is produced. Printed slips kept literal: no printed word, symbol, hypothesis or number of the counted statement or of its proof was altered. Layout and class: the article's own class is \texttt{amsart} with packages including inputenc, lmodern, fontenc, babel, ifpdf, geometry, xcolor, hyperref, amsaddr, graphicx, amsmath, amsthm, amssymb, enumitem; the wrapper is the lane's pinned \texttt{amsart} editorial wrapper at 10pt on A4, which restates the theorem-like environments the retained text needs and adds the article's own macro definitions that the retained text needs (none), with extra wrapper packages (bm, bbm, dsfont, xspace, cancel, mathtools, cleveref). The delivered numbering is the unit's own. Assets: no retained span contains a figure environment, a graphics-inclusion command, a PDF-page inclusion, a table or a TikZ picture; no image file is copied into this package, the package declares no asset dependency and no image file is redistributed with it. Licence: version arXiv:2510.11894v1 of this article is distributed under the Creative Commons Attribution 4.0 International licence, as recorded in the arXiv licence metadata for this exact version and shown on the abstract page https://arxiv.org/abs/2510.11894v1. A full rights notice is delivered at licenses/arxiv-2510.11894-CC-BY-4.0.txt. Characters: every retained excerpt is pure ASCII, so the delivered engine, XeLaTeX with its default Latin Modern text font, reports no missing character.
\begin{lemma}\label{lem:block-diag}
    The Laplacian matrix $\mathbf{L}(G_k)$ admits a block diagonalization of the form
        \begin{align}\label{eq:lap-block-repn}
            \mathbf{L}(G_k) = \widetilde{\mathbf{U}} \begin{pmatrix}
            (\widetilde{\Delta}_k^{(0)}) & & \\
            & (\Delta_k^{(3)}) & \\
            & & (\Delta_k^{(3)})
        \end{pmatrix} \widetilde{\mathbf{U}}^\top,
        \end{align}
    where
        \begin{align*}
            \widetilde{\Delta}_k^{(0)}=
            \begin{pmatrix}
            3 & -\sqrt3 &        &        &        &   \\
            -\sqrt3 & 2 & -1     &        &        &    \\
                    & -1 & 2 & -1&        &        \\
                    &    & \ddots & \ddots & \ddots &   \\
                    &    &        & -1 & 2 & -\sqrt3 \\
                   &    &        &    & -\sqrt3 & 3
            \end{pmatrix}\in\mathbb{R}^{(k+2)\times (k+2)}, 
        \end{align*}
        \begin{align*}
            \Delta_k^{(3)} &= \begin{pmatrix}
            5 & -1 &        &        & \\
            -1& 5  & -1     &        &  \\
               &-1 & 5      & \ddots &  \\
               &   & \ddots & \ddots & -1\\
              &   &        & -1 & 5
            \end{pmatrix}\in\mathbb{R}^{k\times k},
        \end{align*}
    and 
        \begin{align*}
            \widetilde{\mathbf{U}} &= \begin{pmatrix}
                1 & 0\\
                0 & 1
            \end{pmatrix} \oplus \bigoplus_{j=1}^{k} \begin{pmatrix}
                \tfrac{1}{\sqrt{3}} & \tfrac{1}{\sqrt{2}}&\tfrac{1}{\sqrt{6}}\\
                \tfrac{1}{\sqrt{3}} & -\tfrac{1}{\sqrt{2}}& \tfrac{1}{\sqrt{6}}\\
                \tfrac{1}{\sqrt{3}} & 0 & -\tfrac{2}{\sqrt{6}}\\
            \end{pmatrix}.
        \end{align*}
\end{lemma}

\begin{proof}
    We begin by setting, for $0\leq i\leq 2$, the vectors $\mathbf{u}_i\in\mathbb{R}^3$ given by
        \begin{align*}
            \mathbf{u}_0&=\tfrac1{\sqrt3}(1,1,1),&
            \mathbf{u}_1&=\tfrac1{\sqrt2}(1,-1,0),&
            \mathbf{u}_2&=\tfrac1{\sqrt6}(1,1,-2).
        \end{align*}
    These vectors form an orthonormal basis for $\mathbb{R}^3$ and it is straightforward to verify that the cycle Laplacian matrix $\mathbf{L}(C_3)$ satisfies
        \begin{align*}
            \mathbf{L}({C_3})\mathbf{u}_0&=0,& 
            \mathbf{L}({C_3})\mathbf{u}_1&=3\mathbf{u}_1,&
            \mathbf{L}({C_3})\mathbf{u}_2&=3\mathbf{u}_2.
        \end{align*}
    Any function $f:V(C_3)\times V(P_k)\to\mathbb R$ therefore decomposes uniquely as
        \begin{align}\label{eq:expansion-fourier}
        f(i,j)=
            f^{(0)}(j)\mathbf{u}_0(i)+f^{(1)}(j)\mathbf{u}_1(i)+f^{(2)}(j)\mathbf{u}_2(i),
        \end{align}
    where the scalar coefficients $f^{(s)}(j)$, $s=0,1,2$, depend only on
    the index $j$. If we identify the linear space $\ell(V(C_3\times P_k))$ with $\mathbb{R}^{3\times k}$,~\cref{eq:expansion-fourier} can be expressed as
        \begin{align*}
            \begin{pmatrix}
            f(0,0) & f(0,1) & \cdots & f(0,k-1)\\
            f(1,0) & f(1,1) & \cdots & f(1,k-1)\\
            f(2,0) & f(2,1) & \cdots & f(2,k-1)
            \end{pmatrix}
            &=
            U
            \begin{pmatrix}
            f^{(0)}(0) & f^{(0)}(1) & \cdots & f^{(0)}(k-1)\\
            f^{(1)}(0) & f^{(1)}(1) & \cdots & f^{(1)}(k-1)\\
            f^{(2)}(0) & f^{(2)}(1) & \cdots & f^{(2)}(k-1)
            \end{pmatrix}
        \end{align*}
    with $U = \begin{pmatrix}
        \mathbf{u}_0 & \mathbf{u}_1 & \mathbf{u}_2
    \end{pmatrix}\in\mathbb{R}^{3\times 3}$. Let $\mathcal L$ denote the Laplacian of $C_3\times P_k$ {before} the caps $x,y$ are attached. For fixed $0\leq i \leq 2$ and $0\leq j\leq k-1$ the neighbors of $(i,j)$ are
        \begin{align*}
            (i\pm 1,j),\quad (i,j-1),\quad (i,j+1)
        \end{align*}
    with the obvious adjustments when $j = 0, k-1$. Therefore if $f\in \ell(V(C_3\times P_k))$, one has
        {\footnotesize\begin{align*}
            (\mathcal L f)(i,j)= \begin{cases}
                3f(i,j)-\bigl[f(i+1,j)+f(i-1,j)\bigr] - f(i,j+1) &\text{ if } j = 0,\\
                4f(i,j)-\bigl[f(i+1,j)+f(i-1,j)\bigr] -\bigl[f(i,j-1)+f(i,j+1)\bigr] &\text{ if } 0 < j < k-1,\\
                3f(i,j)-\bigl[f(i+1,j)+f(i-1,j)\bigr] - f(i,j-1) &\text{ if } j = k-1.\\
            \end{cases}
        \end{align*}}
    Inserting the expansion set up in~\cref{eq:expansion-fourier}, we have that for $0\leq i\leq 2$ and $0 < j < k-1$ fixed:
        \begin{align*}
            (\mathcal{L}f)(i, j) &= \sum_{s=0}^2 4f^{(s)}(j)\mathbf{u}_s(i)-\bigl[  f^{(s)}(j)\mathbf{u}_s(i+1) + f^{(s)}(j)\mathbf{u}_s(i-1) \bigr]\\
            &\qquad\qquad-\bigl[ f^{(s)}(j-1)\mathbf{u}_s(i) +f^{(s)}(j+1)\mathbf{u}_s(i)\bigr]\\
            &= \sum_{s=0}^2 f^{(s)}(j)(\mathbf{L}(C_3)\mathbf{u}_s)(i) + \mathbf{u}_s(i)(2f^{(s)}(j) -f^{(s)}(j-1) -f^{(s)}(j+1)) \\
            &=\sum_{s=0}^2 \mathbf{u}_s(i)\left[(2+\lambda_s)f^{(s)}(j) -f^{(s)}(j-1) -f^{(s)}(j+1)\right]
        \end{align*}
    where $\lambda_0 = 0$ and $\lambda_1 = \lambda_2= 3$. Similarly, if $j=0, k-1$, we have
        \begin{align*}
            (\mathcal{L}f)(i, j) &= \sum_{s=0}^2 \mathbf{u}_s(i)\left[(2+\lambda_s)f^{(s)}(j) -f^{(s)}(j+1)\right],\quad j=0\\
            (\mathcal{L}f)(i, j) &= \sum_{s=0}^2 \mathbf{u}_s(i)\left[(2+\lambda_s)f^{(s)}(j) -f^{(s)}(j-1)\right],\quad j=k-1.
        \end{align*}
    Inspecting the expressions above, we define the ``longitudinal'' operators $\Delta_k^{(0)}, \Delta_k^{(3)}:\ell(V(P_k))\rightarrow \ell(V(P_k))$ given by their action on functions $g\in \ell(V(P_k))$ as follows:
        \begin{align*}
            (\Delta_k^{(0)}g)(j) &= \begin{cases}
                2g(j)-g(j+1) & \text{ if }j=0,\\
                -g(j-1)+2g(j)-g(j+1) & \text{ if }0<j<k-1,\\
                -g(j-1) + 2g(j) & \text{ if }j=k-1,\\
            \end{cases}\quad g\in\ell(V(P_k)),
        \end{align*}
    and
        \begin{align}\label{eq:anti-block}
            (\Delta_k^{(3)}g)(j) &= \begin{cases}
                5g(j)-g(j+1) & \text{ if }j=0,\\
                -g(j-1)+5g(j)-g(j+1) & \text{ if }0<j<k-1,\\
                -g(j-1) + 5g(j) & \text{ if }j=k-1,\\
            \end{cases}\quad g\in\ell(V(P_k)).
        \end{align}
    We call $\Delta_k^{(0)}$ the {symmetric block} ($\lambda_0=0$), and
    $\Delta_k^{(3)}$ the {antisymmetric block}, and identify the operators with their matrix representations under the standard basis. By~\cref{eq:expansion-fourier} and the preceding, we have that the Laplacian $\mathbf{L}(C_3\times P_k)$ admits the decomposition
        \begin{align*}
            \mathbf{L}(C_3\times P_k) = \mathbf{U} \begin{pmatrix}
            \Delta_k^{(0)} & & \\
            & \Delta_k^{(3)} & \\
            & & \Delta_k^{(3)}
        \end{pmatrix} \mathbf{U}^\top,\qquad \mathbf{U} = \mathrm{diag}(\underbrace{U, U, \dotsc,  U}_{k\text{ times}}).
        \end{align*}
    Now suppose we add to $C_3\times P_k$ the two additional ``cap'' vertices $x, y$ and thereby obtain the graph $G_k$. The space $\ell(V(G_k))$ can be identified as $\mathbb{R}\oplus \mathbb{R} \oplus \ell(V(C_3\times P_k))$. If we set 
        \begin{align*}
            \widetilde{\mathbf{U}} = \mathrm{diag}(I_2, \underbrace{U, U, \dotsc,  U}_{k\text{ times}}),
        \end{align*}
    then we have
        \begin{align}\label{eq:block-decomp-L}
            \mathbf{L}(G_k) = \widetilde{\mathbf{U}} \begin{pmatrix} 
            \mathbf{A}_1 & \mathbf{A}_2 & \mathbf{A}_3 & \mathbf{A}_3\\
            \mathbf{A}_2^\top &\Delta_k^{(0)} & & \\
            \mathbf{A}_3^\top && \Delta_k^{(3)} & \\
            \mathbf{A}_3^\top && & \Delta_k^{(3)}
        \end{pmatrix} \widetilde{\mathbf{U}}^\top
        \end{align}
    where $\mathbf{A}_1\in\mathbb{R}^{2\times 2}$, $\mathbf{A}_2\in\mathbb{R}^{2\times k}$, $\mathbf{A}_3\in\mathbb{R}^{2\times k}$ are to be determined. We claim $\mathbf{A}_3=0$. Let $0\leq j \leq k-1$ be fixed. Let $(\mathbf{A}_3)_{x, j}$ denote the entry of $\mathbf{A}_3$ in the first row (indexed by $x$) and the $j$-th column. We write
        \begin{align*}
            (\mathbf{A}_3)_{x, j} = \mathbf{1}_x^\top \widetilde{\mathbf{U}}^\top \mathbf{L}(G_k)\widetilde{\mathbf{U}} \mathbf{1}_{k+2+j}
        \end{align*}
    Here, $\mathbf{1}_x$ is the indicator vector of the first coordinate (indexed by $x$) and $\mathbf{1}_{k+2+j}$ is the indicator vector of the $(k+2+j)$-th coordinate, which is chosen so as to capture the entry of $\mathbf{A}_3$ as it appears in the first row of the block decomposition of $\mathbf{L}(G_k)$ in~\cref{eq:block-decomp-L}. Note that
        \begin{align*}
            (\widetilde{\mathbf{U}}\mathbf{1}_x)^\top \mathbf{L}(G_k) &= \mathbf{1}_x^\top \mathbf{L}(G_k)\\
            &= 3\mathbf{1}_x^\top - \mathbf{1}_{(0, 0)}^\top - \mathbf{1}_{(1, 0)}^\top - \mathbf{1}_{(2, 0)}^\top
        \end{align*}
    And therefore
        \begin{align*}
            (\mathbf{A}_3)_{x,j}
            &=(3\mathbf{1}_x^\top - \mathbf{1}_{(0,0)}^\top - \mathbf{1}_{(1,0)}^\top - \mathbf{1}_{(2,0)}^\top)\,\widetilde{\mathbf{U}}\,\mathbf{1}_{k+2+j}\\
            &=-\sum_{i=0}^2 \mathbf{1}_{(i,0)}^\top\,\widetilde{\mathbf{U}}\,\mathbf{1}_{k+2+j}.
        \end{align*}
    If the nonzero entries of $\widetilde{\mathbf{U}}\,\mathbf{1}_{k+2+j}$ corresponds to $\mathbf{u}_s$ for $s=1,2$ at level $j=0$, then
        \begin{align*}
            \mathbf{1}_{(i,0)}^\top\,\widetilde{\mathbf{U}}\,\mathbf{1}_{k+2+j}
            = \mathbf{u}_s(i),
            \quad
            \sum_{i=0}^2 \mathbf{u}_s(i)=0,
        \end{align*}
    otherwise the inner product vanishes. In either case the sum is zero, so $(\mathbf{A}_3)_{x,j}=0$.  A similar calculation for the row indexed by $y$ shows $(\mathbf{A}_3)_{y,j}=0$ for all $j$.  Therefore $\mathbf{A}_3=0$ and we have a block diagonalization of $\mathbf{L}(G_k)$ in the form
        \begin{align*}
            \mathbf{L}(G_k) = \widetilde{\mathbf{U}} \begin{pmatrix}
            \widetilde{\Delta}_k^{(0)} & & \\
            & \Delta_k^{(3)} & \\
            & & \Delta_k^{(3)}
        \end{pmatrix} \widetilde{\mathbf{U}}^\top, \qquad \widetilde{\Delta}_k^{(0)} = P\begin{pmatrix}
            \mathbf{A}_1 & \mathbf{A}_2 \\
            \mathbf{A}_2 ^\top & \Delta_k^{(0)}
        \end{pmatrix}P^\top,
        \end{align*}
    where $P$ is the permutation matrix which shifts the coordinate indexing the vertex $y$ to the last slot of the $(k+2)\times (k+2)$ matrix. From~\cref{eq:anti-block}, it follows that
        \begin{align*}
            \Delta_k^{(3)} &= \begin{pmatrix}
            5 & -1 &        &        & \\
            -1& 5  & -1     &        &  \\
            &-1 & 5      & \ddots &  \\
            &   & \ddots & \ddots & -1\\
            &   &        & -1 & 5
            \end{pmatrix}
        \end{align*}
    On the other hand, recall that $\widetilde{\Delta}_k^{(0)} \in\mathbb R^{(k+2)\times(k+2)}$ is obtained from the block
        \begin{align*}
            \begin{pmatrix}
                \mathbf{A}_1 & \mathbf{A}_2\\[2pt]
                \mathbf{A}_2^{\top} & \Delta_k^{(0)}
                \end{pmatrix}
        \end{align*}
    after the permutation $P$ that sends the coordinate indexed by the vertex $y$ to the last position. It therefore suffices to identify the matrices $\mathbf{A}_1$ and $\mathbf{A}_2$ explicitly. For the caps $x$ and $y$ one has
        \begin{align*}
            \mathbf{L}(G_k)\mathbf{1}_x=3\mathbf{1}_x-\sum_{i=0}^2\mathbf{1}_{(i,0)},
            \qquad
            \mathbf{L}(G_k)\mathbf{1}_y=3\mathbf{1}_y-\sum_{i=0}^2\mathbf{1}_{(i,k-1)},
        \end{align*}
    so that, in the $\{x,y\}$-coordinates,
        \begin{align*}
            \mathbf{A}_1=\begin{pmatrix}3&0\\[2pt]0&3\end{pmatrix}.
        \end{align*}
    Now let $0\le j\le k-1$ be fixed and let $\mathbf{v}^{(0)}_j:=\mathbf{U}\mathbf{1}_{3j}$ contain a copy of $\mathbf{u}_0$ at level $j$. Then it holds
        \begin{align*} 
            \mathbf{1}_{(i,0)}^{\top}\mathbf{v}^{(0)}_j=
            \begin{cases}
            \frac1{\sqrt3},&j=0,\\
            0,&j\ne 0,
            \end{cases}
            \qquad
            \mathbf{1}_{(i,k-1)}^{\top}\mathbf{v}^{(0)}_j=
            \begin{cases}
            \frac1{\sqrt3},&j=k-1,\\
            0,&j\ne k-1.
            \end{cases}
        \end{align*}
    Using the expression for $\mathbf{L}(G_k)\mathbf{1}_x$ above, we obtain, 
        \begin{align*}
            (\mathbf{A}_2)_{xj}=\mathbf{1}_x^{\top}\mathbf{L}(G_k)\mathbf{v}^{(0)}_j
            =-\,\sum_{i=0}^2\mathbf{1}_{(i,0)}^{\top}\mathbf{v}^{(0)}_j
            =-\sqrt3\,\delta_{j,0}.
        \end{align*}
    An identical calculation with $\mathbf{1}_y$ gives
        \begin{align*}
            (\mathbf{A}_2)_{yj}=-\sqrt3\,\delta_{j,k-1}.
        \end{align*}
    Hence
        \begin{align*}
            \mathbf{A}_2=\begin{pmatrix}
                -\sqrt3 & 0 &\cdots &0\\
                0&\cdots &0 &-\sqrt3
                \end{pmatrix},
        \end{align*}
    where the first (resp.\ last) column corresponds to $j=0$
    (resp.\ $j=k-1$). Inserting $\mathbf{A}_1$, $\mathbf{A}_2$, and
        \begin{align*}
            \Delta_k^{(0)}=\begin{pmatrix}
                2&-1\\[2pt]-1&\ddots&\ddots\\
                &\ddots&2&-1\\
                &&-1&2
                \end{pmatrix}
        \end{align*}
    into $\begin{pmatrix}\mathbf{A}_1&\mathbf{A}_2\\[2pt]\mathbf{A}_2^{\top}&\Delta_k^{(0)}\end{pmatrix}$, and then apply the permutation $P$ which moves the row/column indexed by $y$ to the bottom-right corner.  The result is the tridiagonal matrix
        \begin{align*}
            \widetilde{\Delta}_k^{(0)}=
            \begin{pmatrix}
            3 & -\sqrt3 &        &        &        & 0  \\
            -\sqrt3 & 2 & -1     &        &        &    \\
                    & -1 & 2 & -1&        &        \\
                    &    & \ddots & \ddots & \ddots &   \\
                    &    &        & -1 & 2 & -\sqrt3 \\
            0       &    &        &    & -\sqrt3 & 3
            \end{pmatrix},
        \end{align*}
    The claim follows.
\end{proof}

\end{document}
